\section{问题分析}

\subsection{问题一的分析}

问题一要求在示向度误差有界的条件下，确定多个检测点交会所得定位区域的直径，并判定以该直径为直径的圆能否覆盖该区域。其难点在于示向度误差在同一地点固定、重复检测无法削弱误差，定位区域由多个半平面求交形成，且以区域直径为直径的圆对区域的覆盖关系不能由直径直接判定。

由于每条示向度读数仅把干扰源约束在一个对称扇形内，我们将该角度约束表示为两个有向半平面约束，构建“半平面交定位模型”，取所有约束之交为定位区域。对非空有界的凸多边形，最远点对必在顶点之间，因此先求半平面交、再枚举顶点对即可得到区域直径，空集、无界及退化情形一并纳入处理。在此基础上，我们比较以该直径为直径的圆与最小包围圆，并通过合法测向约束构造反例，检验该圆覆盖该区域的适用范围。由此得到区域直径计算与覆盖性检验的求解路线。

\subsection{问题二的分析}

问题二要求已知首次示向度，确定第二检测点的选择策略与候选区域，使定位精度达到最优。难点在于第二次检测的反馈有方向、近距离与无信号三类，且两次反馈受同一个未知接收半径约束，不能只依据两条视线的交会角评价选点效果。

为此，我们建立三类反馈下的定位区域更新规则，构建“基于最坏定位误差的第二检测点选址模型”。为了同时给出预测位置与误差尺度，我们以更新区域的最小包围圆圆心作为位置估计、以半径衡量最坏预测误差，并依据最坏情形最优准则，选取最坏反馈误差最小者作为最优检测点。求解时采用粗网格与局部增密，由此完成第二检测点的选点方案。

\subsection{问题三的分析}

问题三要求目标区域内全部为全向源且数量未知时，制定自动搜索、定位与清除的策略与算法，使定位清除总时间最短。其难点在于源的数量与位置均未知、检测范围有限，搜索、定位与清除在一次行进中协同完成。

为此，我们用机器狗位置、当前频道及各频道的检测信息描述系统状态，将移动、切频、检测与清除开销统一换算为时间，构建“基于覆盖约束的搜索-清除联合规划模型”。对于全向源，无信号检测可排除检测点周围的一个圆盘，据此按频道累计排查范围、以覆盖约束保证剩余区域得到检测。在此基础上，我们将搜索、定位与清除共同纳入“滚动规划算法”，每次执行当前路线的首项任务，并根据新反馈更新状态与路线，评价时统计完整退出耗时。由此得到多源情形下的滚动协同搜索清除方案。

\subsection{问题四的分析}

问题四在问题三的基础上引入定向源，检测对象由全向源变为全向源与定向源的混合，要求给出该条件下的检测与清除策略及算法。其难点在于定向源的示向度读数只在特定朝向下可获取，单个无信号反馈的成因有距离过远与检测点位于辐射角之外两类，因而不能直接排除检测点周围的圆盘。

为此，我们沿用问题三的滚动协同框架，利用多个检测点的距离限制与凸包包围关系，构造不依赖具体发射朝向的排查条件，构建“方向覆盖约束下的搜索-清除联合规划模型”。在保留完整排查条件的前提下，我们调整检测点位置、联合安排巡测与已发现源的定位清除，并利用有限的位置、半径与朝向组合估计不同路线的后续检测与定位代价。由此得到混合源情形下的搜索清除方案。
